% =====================================================================
%  Chapter 22 — The \idxw{Baseline} the Moon Gives and the Baseline\idx{Baseline} It Does Not
% =====================================================================

\chapter{The Baseline the Moon Gives and the Baseline It Does Not}
\label{ch:baseline}

Chapter 20 established the requirement: an asteroid shape imager needs a
baseline\idx{Baseline} of hundreds of metres to a few kilometres, matched to the target's
angular size. This chapter asks whether the Moon can provide it, and the
answer is yes for a specific reason that has nothing to do with the Moon
being far away.

\section{The construction problem and the atmosphere problem}

A terrestrial optical array with a baseline\idx{Baseline} of a kilometre has three
simultaneous difficulties, and they are of different kinds.

The \emph{atmospheric} difficulty is that a kilometre of atmosphere is not a
stable optical element. The index of refraction along a kilometre of air
varies with temperature and pressure, and the path length between two
apertures separated by a kilometre changes by a comparable amount. The
optical path difference that an interferometer measures has to be stable to
a small fraction of a wavelength, and a kilometre of air is not.

The \emph{metrological} difficulty is that the two apertures must be
surveyed relative to each other to a fraction of the path stability
required, and a terrestrial baseline is subject to seasonal ground motion,
to thermal expansion of the structures that carry it, and to the slow
deformation of the crust.

The \emph{construction} difficulty is that a kilometre-long structure has to
be built, aligned and maintained on a surface that is not flat, and that
has to be done in a way that survives weather.

The important observation is that these three difficulties are not
independent. On the Earth they are coupled through the atmosphere: the
instrument has to be built on the ground because the atmosphere is above the
ground, and it has to be metrologically stable because the atmosphere is
changing, and it has to be maintained because the ground moves.

On the Moon, the coupling is broken. The ground is rigid and airless and
slowly rotating, with no thermal expansion of a truss to worry about and no
weather and no hydrology. A baseline on the Moon is defined by the distance
between two surveyed anchor points, and the surveying problem is a lunar
geodesy problem rather than an optical-metrology\idx{Metrology} problem.

\begin{seamnote}{This is the defensible lunar claim, and it is narrow}
Not ``the Moon offers a unique vantage point''. Not ``the Moon is closer to
the objects''. The claim is that the lunar surface is the only site where
the kilometre-scale baseline that shape imaging requires can be built,
because the construction problem can be separated from the atmosphere
problem. That is a much smaller claim than the concept literature's, and it
is the one that survives the arithmetic.
\end{seamnote}

\section{The stability argument, and what it is worth}

The concept literature's stability statement is worth reading precisely:
``having the interferometer fixed to the lunar surface keeps the baseline
very stable, reducing measurement uncertainties''.

The physical content is this. A baseline fixed to a rigid, airless body
does not move relative to the body, and the body's orientation relative to
the distant stars is predictable from its rotation. The vector from one
aperture to another therefore has a known direction to a precision set by
the anchor-point survey, and the direction is not corrupted by the
atmosphere's index variations or by the ground's seasonal motion.

What the argument does not say is that the Moon is a good place to build
anything. The regolith is abrasive, the night is long, and the dust is a
mitigation problem with a demonstrated but unproven long-term solution.
Chapter 24 takes the engineering envelope apart. The point of the stability
argument is narrower: it identifies a property that a terrestrial array
cannot have, and that a lunar array can have at a scale that matters.

\section{The baseline the Moon does not give}

Two things have to be said about what the lunar surface does not provide,
and the first of them is the more awkward of the two.

It does not provide the length of baseline that the radio interferometry of
Chapter 8 used. The Very Long Baseline Array reached 1,200 to 2,300
kilometres with the planets as reflectors, and the effective baseline of
that technique came from the Earth's own diameter. A lunar surface array
cannot reach that, and does not need to: the optical wavelength is 750
nanometres rather than 4.2 centimetres, and the required baseline is
correspondingly shorter.

It does not provide proximity. The Moon is not closer to the near-Earth
objects than the Earth is; at a close approach it may be closer or
considerably farther, depending on the geometry. The lunar advantage is
not range. It is the stability of the baseline, the length of the observing
window, and the Sun-ward access.

\begin{table}[!htbp]
\centering
\small
\caption{The three lunar advantages against the terrestrial alternative, and
what each is worth to the shape measurement.}
\label{tab:baseline-compare}
\begin{tabularx}{\textwidth}{@{}P{1.22in}P{1.34in}Y@{}}
\toprule
\rowcolor{panel}
\textbf{Property} & \textbf{Terrestrial array} & \textbf{Lunar array} \\
\midrule
Baseline length & a few hundred metres is already at the limit & a few
hundred metres to kilometres is a construction problem, not an optical one \\
\rowcolor{panel}
Path stability & atmospheric and seasonal & set by the anchor-point survey
alone \\
Observing window & hours, per night & days \\
\rowcolor{panel}
Orthogonal orientation & a second instrument, or a transport delay & 6.83
days, delivered by the rotation \\
Sun-ward access & structurally excluded & available at a polar site \\
\bottomrule
\end{tabularx}
\end{table}

\section{The array the requirement implies}

Putting Chapter 20's requirement together with this chapter's engineering
gives the array specification, and it is worth stating in the form a
programme would use.

The array is a set of fixed apertures at surveyed separations spanning
roughly 300 metres to 3 kilometres, at a site with a low horizon in the
direction of the objects of interest. The station selects the aperture pair
whose separation matches the target's angular size, so that the number of
samples across the diameter is in the range four to seven.

\begin{table}[!htbp]
\centering
\small
\caption{The matched separations the array has to provide, for the class of
interest, at four samples and seven samples across the diameter.}
\label{tab:array-separations}
\begin{tabularx}{\textwidth}{@{}P{0.78in}P{0.90in}P{1.02in}Y@{}}
\toprule
\rowcolor{panel}
\textbf{Diameter} & \textbf{Distance} & \textbf{Four samples} &
\textbf{Seven samples} \\
\midrule
60 m & 0.1 AU & 748 m & 1,309 m \\
\rowcolor{panel}
60 m & 0.2 AU & 1,496 m & 2,618 m \\
60 m & 0.5 AU & 3,740 m & 6,545 m \\
\rowcolor{panel}
140 m & 0.1 AU & 321 m & 561 m \\
140 m & 0.2 AU & 641 m & 1,122 m \\
\rowcolor{panel}
140 m & 0.5 AU & 1,603 m & 2,805 m \\
\bottomrule
\end{tabularx}
\end{table}

The separations in \tabref{tab:array-separations} span a factor of twelve,
from 321 metres to 6,545 metres, and they change with the distance to the
target, which changes with the object's orbit and therefore with the season.
A single-baseline instrument cannot serve this table. An array with a
discrete set of anchor separations can, and the selection rule that chooses
the pair is a rule about the target, not about the instrument.

\begin{keytakeaways}
\item A terrestrial kilometre-scale array is simultaneously an atmospheric, a
metrological and a construction problem; on the Moon the construction
problem separates from the atmosphere problem.
\item The Moon does not give a longer baseline than the radio technique of
Chapter 8, and it does not give proximity; it gives stability, window length
and Sun-ward access.
\item The baseline orientation requirement is met by the Moon's own rotation
in 6.83 days, with no moving hardware.
\item The array must offer separations from 321 metres to 6,545 metres, which
means selectable anchor pairs rather than a single interferometer.
\end{keytakeaways}
